%% GENERATED by tools/gen_latex.py -- edit content/ch19.py instead.
%% Build:  latexmk -pdf ch19.tex     (or: pdflatex ch19.tex, twice)
\documentclass[aspectratio=169,11pt,t]{beamer}
\usepackage{itbpro}

\coursetitle{The Future of AI}
\coursesubtitle{Where deep learning falls short, why scale does not fix it, and what a system would need in order to handle a situation nobody anticipated.}
\coursekicker{Chapter 19}
\coursesource{Chollet \& Watson, Deep Learning with Python 3e --- chapter 19}
\courseduration{2.5 hours (2 sessions)}
\coursepresenter{Prof. Bambang Riyanto Trilaksono}{Lead Instructor}
\coursebrand{AI for Professional \textperiodcentered\ ITB \textperiodcentered\ Ch 19}

\title{The Future of AI}
\author{Prof. Bambang Riyanto Trilaksono}
\date{}

\begin{document}
\covertitle

\framekicker{Learning outcomes}
\begin{frame}[shrink=25]{Session Objectives}
\leadpar{By the end of this session, participants can:}
\begin{isteps}
\item State the \textbf{four structural limitations} of deep learning, and give a concrete failure for each.
\item Explain why an LLM's success on a problem tracks \textbf{familiarity, not complexity}.
\item Explain what \textbf{prompt engineering actually is} in terms of the interpolative database, and what its necessity implies.
\item Argue why \textbf{scaling laws do not answer} any of these limitations.
\item Distinguish \textbf{local, broad, and extreme generalization}, and place current systems on that spectrum.
\item Define intelligence as an \textbf{efficiency ratio} between information available and future operating area.
\item Explain what \textbf{ARC-AGI} measures, why it resisted a 50,000$\times$ scale-up, and what test-time adaptation changed in 2024.
\item Compare the \textbf{two poles of abstraction} and say why a complete system needs both.
\end{isteps}

\vskip 6pt
\vskip 4pt
\reslink{Book}{Chapter 19 --- full text}{https://deeplearningwithpython.io/chapters/chapter19\_future\_of\_ai}
\vskip 2pt
\end{frame}

\coursesection{01}{The limitations of deep learning}{To use a tool well, be aware of its limitations, not just its strengths.}

\framekicker{Section 19.1.1}
\begin{frame}[shrink=25]{A parametric curve is a database, and a static one}
{\fontsize{9.4}{12.6}\selectfont\color{fg2}Deep learning models are \textbf{big parametric curves fitted to large datasets}. That is the source of their power --- easy to train, scaling well in both model and dataset size. It is also the source of significant weaknesses.\par}\vskip 4pt
\begin{center}
\includegraphics[width=0.92\linewidth,height=0.60\textheight,keepaspectratio]{figs/ch19-static-db.tex.pdf}
\end{center}
\vskip 3pt{\color{ink3}\fontsize{7.4}{9.4}\selectfont Chapter 15 called it an interpolative database. The word that matters here is \textit{static}.\par}
{\fontsize{9.4}{12.6}\selectfont\color{fg2}\textbf{The only thing you can do with a static database is information retrieval.} At inference time you had better hope the situations the model faces are inside the training distribution.\par}\vskip 4pt
\end{frame}

\framekicker{Section 19.1.1}
\begin{frame}[shrink=25]{The leopard-print sofa}
\leadpar{A model trained on ImageNet will classify a \textbf{leopard-print sofa} as an actual leopard. Sofas were not part of its training data.}
{\fontsize{9.4}{12.6}\selectfont\color{fg2}This applies equally to the largest generative models. It is frequently claimed that LLMs perform \textbf{in-context learning} --- picking up new skills from a few examples. There is overwhelming evidence that what they are actually doing is \textbf{fetching vector functions memorized during training} and reapplying them.\par}\vskip 4pt
{\fontsize{9.4}{12.6}\selectfont\color{fg2}By learning next-token prediction across a web-sized dataset, an LLM has collected millions of potentially useful mini text-processing programs, and can be prompted into reusing them. \hilite{Show it something with no direct equivalent in its training data, and it is helpless.}\par}\vskip 4pt
\end{frame}

\framekicker{Section 19.1.1 $\cdot$ the diagnosis}
\begin{frame}[shrink=25]{Familiarity, not complexity}
\begin{center}
\includegraphics[width=0.92\linewidth,height=0.60\textheight,keepaspectratio]{figs/ch19-familiarity.tex.pdf}
\end{center}
\vskip 3pt{\color{ink3}\fontsize{7.4}{9.4}\selectfont Figure 19.1's puzzle is easy for you and unsolved by every state-of-the-art model.\par}
{\fontsize{9.4}{12.6}\selectfont\color{fg2}This is the single most useful reframing in the chapter for anyone deciding where to deploy an LLM. \textbf{The question is never "is this task hard?" It is "is this task familiar?"}\par}\vskip 4pt
\end{frame}
\note{Ask the room for a task from their own work that is trivially easy and almost certainly absent from web text. That is the failure case they will hit in production.}

\framekicker{Section 19.1.1 $\cdot$ two well-known failures}
\begin{frame}[shrink=25]{Ten kilos of steel, and Monty Hall}
\begin{iquote}
\quotetext{\textit{"What's heavier, \textbf{10 kilos} of steel or one kilo of feathers?"} --- For months after ChatGPT's release, the answer was that \textbf{they weigh the same}.}\quotecite{Section 19.1.1}
\end{iquote}
{\fontsize{9.4}{12.6}\selectfont\color{fg2}The question \textit{"one kilo of steel or one kilo of feathers?"} appears many times online as a trick question, with that exact answer. The model repeated the memorized answer \textbf{without attending to the actual numbers} or what the query meant.\par}\vskip 4pt
{\fontsize{9.4}{12.6}\selectfont\color{fg2}The same happens with variations of the \textbf{Monty Hall problem}: the canonical answer comes out regardless of whether it makes sense in the altered context.\par}\vskip 4pt
\end{frame}

\framekicker{Section 19.1.1 $\cdot$ how these get fixed}
\begin{frame}[shrink=25]{Twenty-five thousand people playing whack-a-mole}
{\fontsize{9.4}{12.6}\selectfont\color{fg2}Those specific prompts were patched later --- by \textbf{special-casing them}.\par}\vskip 4pt
\begin{minipage}[t]{0.494\textwidth}
\istat{25,000+}{people employed full time providing LLM training data}
\end{minipage}%
\hfill
\begin{minipage}[t]{0.494\textwidth}
\istat{one at a time}{how failing prompts get patched}
\end{minipage}%
\begin{iband}[rose]
LLM maintenance is a constant game of whack-a-mole where failing prompts are patched individually, \textbf{without addressing the underlying issue}. \hilite{Even already-patched prompts will fail again if you make small changes to them.}
\end{iband}
\end{frame}

\framekicker{Section 19.1.2 $\cdot$ figure 19.3}
\begin{frame}[shrink=25]{Adversarial examples: a panda plus a gibbon gradient}
{\fontsize{9.4}{12.6}\selectfont\color{fg2}Chapter 10 showed gradient ascent in input space to generate inputs maximizing a ConvNet filter's activation. The same technique can \textbf{slightly modify an image to maximize the prediction for a chosen class}.\par}\vskip 4pt
{\fontsize{9.4}{12.6}\selectfont\color{fg2}Take a picture of a panda, add a gibbon class gradient, and the network classifies it as a gibbon. The change is \textbf{imperceptible to a human}.\par}\vskip 4pt
\begin{iband}[signal]
This evidences both the brittleness of these models and the \hilite{deep difference between their input-to-output mapping and human perception}. The two are not approximations of each other.
\end{iband}
\end{frame}

\framekicker{Section 19.1.2 $\cdot$ the text equivalent}
\begin{frame}[shrink=25]{Alice has N brothers and M sisters}
\begin{iquote}
\quotetext{\textit{"Alice has N brothers and she also has M sisters. How many sisters does Alice's brother have?"}}\quotecite{The Alice in Wonderland riddle --- Nezhurina et al., arXiv 2406.02061}
\end{iquote}
{\fontsize{9.4}{12.6}\selectfont\color{fg2}The answer is \textbf{M + 1} --- Alice's sisters plus Alice herself. With values commonly found in online instances of the riddle (N = 3, M = 2) an LLM generally answers correctly. \textbf{Tweak the values and you quickly get wrong answers.}\par}\vskip 4pt
{\fontsize{9.4}{12.6}\selectfont\color{fg2}Changing place names, people's names in a paragraph, or variable names in a block of code can \textbf{significantly degrade} performance. None of those changes alters the problem.\par}\vskip 4pt
\end{frame}

\framekicker{Section 19.1.2 $\cdot$ what prompt engineering really is}
\begin{frame}[shrink=25]{Two framings of the same fact}
\begin{minipage}[t]{0.494\textwidth}
\begin{icard}
\cardh{The optimistic framing}
\textit{Your models are better than you know! You just need to use them right.} Adding \textbf{"Please think step by step"} can significantly boost performance on reasoning tasks.
\end{icard}
\end{minipage}%
\hfill
\begin{minipage}[t]{0.494\textwidth}
\begin{icardwarn}
\cardh{The negative framing}
For any query that seems to work, there is a range of \textbf{minor changes that can tank performance}. To what extent does a model understand something if a rewording breaks that understanding?
\end{icardwarn}
\end{minipage}%
\begin{center}
\includegraphics[width=0.92\linewidth,height=0.60\textheight,keepaspectratio]{figs/ch19-prompt-lookup.tex.pdf}
\end{center}
\vskip 3pt{\color{ink3}\fontsize{7.4}{9.4}\selectfont A prompt is an address, not an instruction.\par}
\end{frame}

\framekicker{Section 19.1.2}
\begin{frame}[shrink=25]{Prompt engineering is search through latent space}
{\fontsize{9.4}{12.6}\selectfont\color{fg2}Because the knowledge and programs indexed by the LLM are \textbf{interpolative}, you can move around in latent space to nearby locations. \textit{"Explain Python list sorting, but answer like a buccaneer"} points somewhere very similar --- close, but not identical.\par}\vskip 4pt
{\fontsize{9.4}{12.6}\selectfont\color{fg2}There are thousands of variations, each giving a similar yet slightly different answer. \textbf{There is no a priori reason for your first naive prompt to be optimal.}\par}\vskip 4pt
\begin{iband}[signal]
Prompt engineering is trial-and-error search for the lookup query that performs best. \textbf{It is no different from trying different keywords in a Google search.} If LLMs actually understood what you asked, \hilite{there would be no need for this search process at all}
\end{iband}
\end{frame}
\note{The logical argument here is the sharp one: the information conveyed about your task does not change whether you write 'rewrite' or 'rephrase'. If the output changes, something other than understanding is happening.}

\framekicker{Section 19.1.3}
\begin{frame}[shrink=25]{Memorized programs do not generalize}
{\fontsize{9.4}{12.6}\selectfont\color{fg2}Suppose you only need a well-known program and you know exactly how to address it in latent space. \textbf{You still have a problem}: the programs deep learning models memorize often do not generalize. They work for some inputs and fail for others.\par}\vskip 4pt
{\fontsize{9.4}{12.6}\selectfont\color{fg2}This is especially true of programs encoding \textbf{discrete logic}. Train a Transformer on hundreds of thousands of digit-addition pairs like \inlinecode{"4 3 5 7 + 8 9 3 6"} and you reach very high accuracy --- \textbf{very high, but not 100\%}.\par}\vskip 4pt
\begin{minipage}[t]{0.494\textwidth}
\istat{\textasciitilde{}70\%}{state-of-the-art LLM accuracy on digit addition}
\end{minipage}%
\hfill
\begin{minipage}[t]{0.494\textwidth}
\istat{which digits}{accuracy depends on it --- common digits do better}
\end{minipage}%
{\fontsize{9.4}{12.6}\selectfont\color{fg2}Unless, that is, the model was explicitly hardcoded to execute \inlinecode{4357 + 8936} in Python.\par}\vskip 4pt
\end{frame}

\framekicker{Section 19.1.3 $\cdot$ why}
\begin{frame}[shrink=25]{Most programs cannot be expressed as a deep learning model}
{\fontsize{9.4}{12.6}\selectfont\color{fg2}A deep learning model is a \textbf{static chain of simple, continuous geometric transformations} mapping one vector space into another. That is a good fit for perceptual pattern recognition and a very poor fit for step-by-step discrete logic --- concepts like \textbf{place value} or \textbf{carrying over}.\par}\vskip 4pt
\begin{iband}[signal]
A deep learning model can be interpreted as a kind of program. But inversely, \textbf{most programs cannot be expressed as deep learning models.} For most tasks either no network of reasonable size solves it, or one exists but is not learnable --- the transform is too complex, or the data does not exist.
\end{iband}
\end{frame}

\framekicker{Section 19.1.4 $\cdot$ figure 19.4}
\begin{frame}[shrink=25]{A dim image in a mirror}
{\fontsize{9.4}{12.6}\selectfont\color{fg2}Our understanding of images, sounds, and language is grounded in \textbf{sensorimotor experience}. Models have no access to that, and so cannot understand their inputs in a human-relatable way.\par}\vskip 4pt
\begin{center}
\includegraphics[width=0.92\linewidth,height=0.60\textheight,keepaspectratio]{figs/ch19-mirror.tex.pdf}
\end{center}
\vskip 3pt{\color{ink3}\fontsize{7.4}{9.4}\selectfont What the model learns matches the training data --- not the mental model the data came from.\par}
{\fontsize{9.4}{12.6}\selectfont\color{fg2}\textbf{The models you create will take any shortcut available to fit their training data.}\par}\vskip 4pt
\end{frame}

\framekicker{Section 19.1.4 $\cdot$ the risk}
\begin{frame}[shrink=25]{Theory of mind, applied where it does not belong}
{\fontsize{9.4}{12.6}\selectfont\color{fg2}A fundamental feature of humans is our \textbf{theory of mind} --- our tendency to project intentions, beliefs, and knowledge onto the things around us. Draw a smiley face on a rock and it becomes \textit{happy}, in our minds.\par}\vskip 4pt
{\fontsize{9.4}{12.6}\selectfont\color{fg2}Applied to deep learning, this leads us to believe a model that uses language \textbf{understands} the word sequences it generates the way we do. Then we are surprised when a slight departure from the training patterns produces something absurd.\par}\vskip 4pt
\begin{iband}[rose]
Never fall into the trap of believing neural networks understand the task they perform. \textbf{They were trained on a different, far narrower task than the one you wanted to teach them}: mapping training inputs to training targets, point by point.
\end{iband}
\end{frame}

\coursesection{02}{Scale isn't all you need}{Four or five orders of magnitude larger, and the same problems.}

\framekicker{Section 19.2}
\begin{frame}[shrink=25]{The prevailing narrative, and what it missed}
{\fontsize{9.4}{12.6}\selectfont\color{fg2}In early 2023, at peak LLM hype, GPT-4 had just been released --- essentially a scaled-up GPT-3. Its improved performance seemed to suggest you could just keep going, and that AGI would spontaneously emerge from a GPT-5.\par}\vskip 4pt
{\fontsize{9.4}{12.6}\selectfont\color{fg2}Proponents pointed to \textbf{scaling laws}: an empirical relationship between model and dataset size and performance on specific tasks, suggesting that size reliably and predictably buys performance.\par}\vskip 4pt
\begin{iband}[signal]
The key thing scaling-law enthusiasts miss is that the benchmarks measuring "performance" are effectively \textbf{memorization tests} --- the kind we give university students. LLMs do well by memorizing the answers, and \hilite{cramming more questions and answers in improves the score accordingly.}
\end{iband}
\end{frame}

\framekicker{Section 19.2 $\cdot$ the evidence}
\begin{frame}[shrink=25]{Seven years, four orders of magnitude, same problems}
\leadpar{Scaling has produced \textbf{no progress} on inability to adapt to novelty, oversensitivity to phrasing, or the inability to infer generalizable programs --- because these issues are inherent to \textbf{curve fitting}.}
{\fontsize{9.4}{12.6}\selectfont\color{fg2}Chollet began pointing out these problems in \textbf{2017}. We are still struggling with them, with models four or five orders of magnitude larger and more knowledgeable.\par}\vskip 4pt
\begin{iband}[signal]
\textbf{We have made no progress because the models are still the same.} They have been the same for over seven years: parametric curves fitted to a dataset by gradient descent, using the Transformer architecture.
\end{iband}
\end{frame}

\framekicker{Section 19.2 $\cdot$ stated plainly}
\begin{frame}[shrink=25]{Three things stacking more layers will not solve}
\begin{isteps}
\item Models are \textbf{limited to interpolative programs memorized at training time}. They cannot, on their own, synthesize brand-new programs at inference time to adapt to substantially novel situations.
\item Even within known situations, those interpolative programs suffer \textbf{generalization issues} --- oversensitivity to phrasing and confounder features.
\item Models are \textbf{limited in what they can represent}. Most programs you would wish to learn cannot be expressed as a continuous geometric morphing of a data manifold --- algorithmic reasoning tasks in particular.
\end{isteps}
\end{frame}

\framekicker{Section 19.2.1}
\begin{frame}[shrink=25]{Why brains appeared, and what they were at first}
\begin{center}
\includegraphics[width=0.92\linewidth,height=0.60\textheight,keepaspectratio]{figs/ch19-brain.tex.pdf}
\end{center}
\vskip 3pt{\color{ink3}\fontsize{7.4}{9.4}\selectfont Brains appeared more than half a billion years ago as a way to store and execute behavioural programs.\par}
{\fontsize{9.4}{12.6}\selectfont\color{fg2}Because the source code was \textbf{DNA}, decoded as neural connectivity, evolution could suddenly search over behaviour space in a largely unbounded way. Eyes, ears, mandibles, four legs, twenty-four legs --- \textbf{brains handle any modality you throw at them}.\par}\vskip 4pt
\end{frame}

\framekicker{Section 19.2.2}
\begin{frame}[shrink=25]{Automatons and intelligent agents}
{\fontsize{9.4}{12.6}\selectfont\color{fg2}The field of AI has long suffered from \textbf{conflating intelligence with automation}.\par}\vskip 4pt
\begin{minipage}[t]{0.494\textwidth}
\begin{icard}
\cardh{An automaton}
Static, crafted to accomplish specific things in a specific context. Expose it to something outside what it was programmed for --- by hand, by evolution, or by fitting on a dataset --- and \textbf{it fails}.
\end{icard}
\end{minipage}%
\hfill
\begin{minipage}[t]{0.494\textwidth}
\begin{icardaccent}
\cardh{An intelligent agent}
Adapts \textbf{on the fly} to novel, unexpected situations, using fluid intelligence to find a way forward.
\end{icardaccent}
\end{minipage}%
\begin{iband}[signal]
How do you tell a student who memorized three years of past exam questions from one who understands the subject? \hilite{You give them a brand-new problem.}
\end{iband}
\end{frame}

\framekicker{Section 19.2.2 $\cdot$ figure 19.6}
\begin{frame}[shrink=25]{The spectrum of generalization}
\begin{center}
\includegraphics[width=0.92\linewidth,height=0.60\textheight,keepaspectratio]{figs/ch19-generalization.tex.pdf}
\end{center}
\vskip 3pt{\color{ink3}\fontsize{7.4}{9.4}\selectfont Deep nets generalize to \textit{known unknowns} --- factors of variation anticipated during development and featured in the training data.\par}
{\fontsize{9.4}{12.6}\selectfont\color{fg2}Deep nets generalize \textbf{by interpolation on a manifold}, so any factor of variation in the input space must be captured by the manifold they learn. That is why basic \textbf{data augmentation} helps so much --- and why nothing helps when the variation was never anticipated.\par}\vskip 4pt
\end{frame}

\framekicker{Section 19.2.2 $\cdot$ two thought experiments}
\begin{frame}[shrink=25]{Landing a rocket, and crossing a road}
\begin{minipage}[t]{0.494\textwidth}
\begin{icardwarn}
\cardh{The rocket}
A deep net needs \textbf{tens of thousands or millions of launch trials} --- a dense sampling of the input space. Humans build a physical model, \textit{rocket science}, and land it in one or a few tries.
\end{icardwarn}
\end{minipage}%
\hfill
\begin{minipage}[t]{0.494\textwidth}
\begin{icardwarn}
\cardh{The road}
A net controlling a body would have to \textbf{die many thousands of times} to infer that cars are dangerous. Dropped into a new city, it would relearn most of what it knows.
\end{icardwarn}
\end{minipage}%
{\fontsize{9.4}{12.6}\selectfont\color{fg2}Humans learn safe behaviours \textbf{without dying even once} --- thanks to the power of abstract modelling of novel situations.\par}\vskip 4pt
\end{frame}

\framekicker{Section 19.2.3}
\begin{frame}[shrink=25]{What intelligence is for}
{\fontsize{9.4}{12.6}\selectfont\color{fg2}If the set of situations an organism faced were static and known, behaviour generation would be easy: evolution would find the right behaviours by trial and error and hardcode them into DNA. \textbf{Brains as automatons would already be optimal.}\par}\vskip 4pt
{\fontsize{9.4}{12.6}\selectfont\color{fg2}But as organism and environmental complexity rose, situations became dynamic and unpredictable. \textbf{A day in your life is unlike any day you have experienced, and unlike any day experienced by any of your evolutionary ancestors.}\par}\vskip 4pt
\begin{iquote}
\quotetext{Intelligence is the ability to \textbf{efficiently use the information at your disposal to produce successful behaviour in the face of an uncertain, ever-changing future.}}\quotecite{Section 19.2.3}
\end{iquote}
\end{frame}

\framekicker{Section 19.2.4}
\begin{frame}[shrink=25]{Seventy years of AI, on the same spectrum}
\begin{itable}{>{\raggedright\arraybackslash}p{0.1692\textwidth}>{\raggedright\arraybackslash}p{0.3948\textwidth}>{\raggedright\arraybackslash}p{0.376\textwidth}}
\thead{Era} & \thead{System} & \thead{Where on the spectrum} \\ \midrule
1960s-70s & ELIZA; SHRDLU manipulating objects from language commands & \textbf{Pure automatons} \\
1990s-2000s & Machine learning systems handling some uncertainty & \textbf{Local generalization} \\
2010s & Deep learning: larger datasets, more expressive models & \textbf{Local generalization, expanded} \\
Today & Self-driving cars; a robot that could make coffee in a random kitchen & \textbf{Toward broad generalization} \\
\end{itable}
{\fontsize{9.4}{12.6}\selectfont\color{fg2}The last row is the \textbf{Woz test of intelligence} --- enter a random kitchen and make a cup of coffee. Visible progress is being made, by \textbf{combining deep learning with painstakingly handcrafted abstract models of the world}.\par}\vskip 4pt
\end{frame}

\framekicker{Section 19.2.4 $\cdot$ a proposed renaming}
\begin{frame}[shrink=25]{Artificial cognition, not artificial intelligence}
\begin{iquote}
\quotetext{The "intelligence" label in "artificial intelligence" has been a \textbf{category error}. It would be more accurate to call our field \textit{artificial cognition}, with \textit{cognitive automation} and \textit{artificial intelligence} being two nearly independent subfields within it.}\quotecite{Section 19.2.4}
\end{iquote}
{\fontsize{9.4}{12.6}\selectfont\color{fg2}In that subdivision, AI would be \textbf{a greenfield where almost everything remains to be discovered}.\par}\vskip 4pt
\begin{iband}[signal]
This is not meant to diminish deep learning. Cognitive automation is incredibly useful, and automating tasks from exposure to data alone is \textbf{far more practical and versatile than explicit programming}. \hilite{Doing it well is a game changer for essentially every industry.}
\end{iband}
\end{frame}
\note{Land this carefully for a professional audience. The chapter is not saying the technology is oversold as a product. It is saying the word 'intelligence' is oversold as a description.}

\coursesection{03}{How to build intelligence}{Abstraction acquisition, on-the-fly recombination, and a benchmark designed to resist memorization.}

\framekicker{Section 19.3.1}
\begin{frame}[shrink=25]{The kaleidoscope hypothesis}
{\fontsize{9.4}{12.6}\selectfont\color{fg2}If the future you faced were \textbf{truly} novel --- sharing no common ground with anything you had seen --- you would be unable to react, no matter how intelligent. \textbf{Intelligence works because nothing is ever truly without precedent.}\par}\vskip 4pt
{\fontsize{9.4}{12.6}\selectfont\color{fg2}A person from the 17th century seeing a jet plane might describe \textit{a large, loud metal bird that doesn't flap its wings}. A car is a \textbf{horseless carriage}. Electricity is like water in a pipe; spacetime is like a rubber sheet distorted by heavy objects.\par}\vskip 4pt
\begin{center}
\includegraphics[width=0.92\linewidth,height=0.60\textheight,keepaspectratio]{figs/ch19-kaleidoscope.tex.pdf}
\end{center}
\vskip 3pt{\color{ink3}\fontsize{7.4}{9.4}\selectfont A few beads of coloured glass, reflected by mirrors, produce endless patterns.\par}
\end{frame}

\framekicker{Section 19.3.1}
\begin{frame}[shrink=25]{The analogies are not only in our heads}
{\fontsize{9.4}{12.6}\selectfont\color{fg2}Besides clear-cut explicit analogies, we make \textbf{smaller implicit ones every second, with every thought}. A new supermarket relates to similar stores. A new person reminds you of a few you have met. Even cloud shapes evoke an elephant, a ship, a fish.\par}\vskip 4pt
{\fontsize{9.4}{12.6}\selectfont\color{fg2}And \textbf{physical reality itself is full of isomorphisms}. Electromagnetism is analogous to gravity. Animals are structurally similar because of shared origins. Silica crystals resemble ice crystals.\par}\vskip 4pt
\begin{iband}[signal]
The number of unique \textbf{atoms of meaning} needed to describe the universe you live in is relatively small, and everything around you is a \hilite{recombination of them}
\end{iband}
\end{frame}

\framekicker{Section 19.3.2}
\begin{frame}[shrink=25]{Intelligence in two parts}
\begin{center}
\includegraphics[width=0.92\linewidth,height=0.60\textheight,keepaspectratio]{figs/ch19-two-parts.tex.pdf}
\end{center}
\vskip 3pt{\color{ink3}\fontsize{7.4}{9.4}\selectfont Extract the beads; recombine them on the fly into a model of the situation in front of you.\par}
\begin{iband}[signal]
\textbf{The emphasis on efficiency is crucial.} If you need hundreds of thousands of hours of practice to acquire a skill, you are not very intelligent. If you need to enumerate every possible move on a chessboard to find the best one, you are not very intelligent.
\end{iband}
\end{frame}

\framekicker{Section 19.3.2 $\cdot$ the diagnosis}
\begin{frame}[shrink=25]{Deep learning has one half, and it is the wrong half}
\begin{minipage}[t]{0.494\textwidth}
\begin{icardbad}
\cardh{No on-the-fly recombination}
Models do a decent job of \textbf{acquiring} abstractions at training time via gradient descent, but have \textbf{zero ability to recombine} what they know at test time. They are a static database limited to retrieval --- missing half the picture, arguably the most important half.
\end{icardbad}
\end{minipage}%
\hfill
\begin{minipage}[t]{0.494\textwidth}
\begin{icardbad}
\cardh{Terribly inefficient}
Gradient descent requires vast amounts of data to distill neat abstractions --- \textbf{many orders of magnitude more than humans}.
\end{icardbad}
\end{minipage}%
\end{frame}

\framekicker{Section 19.3.3 $\cdot$ the shortcut rule}
\begin{frame}[shrink=25]{You get what you measure, and nothing else}
\leadpar{\textbf{The shortcut rule}: optimize one success metric and you will achieve your goal --- at the expense of everything in the system your metric did not cover.}
{\fontsize{9.4}{12.6}\selectfont\color{fg2}In 2009 Netflix ran a \$1 million challenge for the best movie recommendation score. \textbf{It never used the winning system} --- far too complex and compute intensive. The winners had optimized for prediction accuracy alone, at the expense of inference cost, maintainability, and explainability.\par}\vskip 4pt
{\fontsize{9.4}{12.6}\selectfont\color{fg2}The same holds in most Kaggle competitions: winning models can \textbf{rarely, if ever, be used in production}. \hilite{Your creations are shaped by the incentives you give yourself.}\par}\vskip 4pt
\end{frame}

\framekicker{Section 19.3.3 $\cdot$ the canonical example}
\begin{frame}[shrink=25]{Deep Blue, and what it taught us about the mind}
{\fontsize{9.4}{12.6}\selectfont\color{fg2}In the 1970s Allen Newell, concerned his field was making no progress toward a theory of cognition, proposed a new grand goal: \textbf{chess playing}. Chess in humans seemed to require perception, reasoning, analysis, memory, study from books. Surely a chess machine would need those too?\par}\vskip 4pt
{\fontsize{9.4}{12.6}\selectfont\color{fg2}In 1997 Deep Blue beat Kasparov. Researchers then had to contend with the fact that it had \textbf{taught them little about human intelligence}. The A-star algorithm at its heart was not a model of the brain and could not generalize beyond similar board games.\par}\vskip 4pt
\begin{iband}[signal]
It turned out to be easier to build an AI that could only play chess than to build an artificial mind --- \textbf{so that is the shortcut researchers took}.
\end{iband}
\end{frame}

\framekicker{Section 19.3.3 $\cdot$ the structural argument}
\begin{frame}[shrink=25]{Fixing the task removes the need for intelligence}
\begin{center}
\includegraphics[width=0.92\linewidth,height=0.60\textheight,keepaspectratio]{figs/ch19-shortcut.tex.pdf}
\end{center}
\vskip 3pt{\color{ink3}\fontsize{7.4}{9.4}\selectfont The driving success metric of AI has been to solve specific tasks --- from chess to Go, MNIST to ImageNet, high school maths to the bar exam.\par}
{\fontsize{9.4}{12.6}\selectfont\color{fg2}Given near-infinite training data, even nearest-neighbour search can play video games with superhuman skill. Given near-infinite hand-written if-then-else statements, so can those --- \textbf{until you make a small change to the rules, the kind a human adapts to instantly}.\par}\vskip 4pt
\end{frame}

\framekicker{Section 19.3.3}
\begin{frame}[shrink=25]{There is no path from skills to intelligence}
\leadpar{Humans can use general intelligence to acquire skill at any new task. \textbf{In reverse, there is no path from a collection of task-specific skills to general intelligence.}}
{\fontsize{9.4}{12.6}\selectfont\color{fg2}By fixing the task you make it possible to describe precisely what needs to be done, and to \textbf{buy} more skill by adding data or hardcoded knowledge --- without increasing generalization power at all.\par}\vskip 4pt
\begin{iband}[signal]
Human-like intelligence is not characterized by skill at any particular task. It is \textbf{the ability to adapt to novelty, to efficiently acquire new skills and master never-before-seen tasks}.
\end{iband}
\end{frame}

\framekicker{Section 19.3.4}
\begin{frame}[shrink=25]{Intelligence as an efficiency ratio}
\begin{center}
\includegraphics[width=0.92\linewidth,height=0.60\textheight,keepaspectratio]{figs/ch19-ratio.tex.pdf}
\end{center}
\vskip 3pt{\color{ink3}\fontsize{7.4}{9.4}\selectfont Fix the information available; measure performance on situations known to be sufficiently different from it.\par}
\begin{iband}[amber]
To avoid cheating you must test only on tasks the system was not programmed or trained to handle --- in fact, \textbf{tasks its creators could not have anticipated}.
\end{iband}
\end{frame}

\framekicker{Section 19.3.4 $\cdot$ figure 19.9}
\begin{frame}[shrink=25]{ARC-AGI}
{\fontsize{9.4}{12.6}\selectfont\color{fg2}In 2018-19 Chollet built the \textbf{Abstraction \& Reasoning Corpus for Artificial General Intelligence}, a benchmark designed to capture that definition. It is approachable by both machines and humans, and looks like a human IQ test --- Raven's progressive matrices.\par}\vskip 4pt
{\fontsize{9.4}{12.6}\selectfont\color{fg2}Each task is explained by three or four \textbf{input grid / output grid} examples. You are given a new input grid and have three tries to produce the correct output before moving on.\par}\vskip 4pt
\begin{minipage}[t]{0.494\textwidth}
\begin{icardaccent}
\cardh{Only unseen tasks}
\textbf{A game you cannot practise for.} Each task has its own unique logic to be understood on the fly; memorizing strategies from past tasks does not transfer.
\end{icardaccent}
\end{minipage}%
\hfill
\begin{minipage}[t]{0.494\textwidth}
\begin{icardaccent}
\cardh{Controlled priors}
All test takers start from \textbf{Core Knowledge priors} --- the knowledge systems humans are born with. Unlike an IQ test, tasks never involve acquired knowledge such as English sentences.
\end{icardaccent}
\end{minipage}%
\end{frame}

\framekicker{Section 19.3.5}
\begin{frame}[shrink=25]{A 50,000$\times$ scale-up bought ten percentage points}
{\fontsize{9.4}{12.6}\selectfont\color{fg2}In 2024 Chollet and Mike Knoop founded the nonprofit \textbf{ARC Prize Foundation}, running a yearly competition with over \$1 million in prizes.\par}\vskip 4pt
\begin{minipage}[t]{0.3253\textwidth}
\istat{\textasciitilde{}50,000$\times$}{base LLM scale-up, GPT-2 (2019) to GPT-4.5 (2025)}
\end{minipage}%
\hfill
\begin{minipage}[t]{0.3253\textwidth}
\istat{0\% $\rightarrow$ \textasciitilde{}10\%}{their ARC-AGI score over that period}
\end{minipage}%
\hfill
\begin{minipage}[t]{0.3253\textwidth}
\istat{> 95\%}{what you, the reader, would score}
\end{minipage}%
\begin{iband}[amber]
Most benchmarks saturated quickly in the age of LLMs, because they can be \textbf{hacked via memorization}. ARC-AGI was designed to resist it. \hilite{Scale up 50,000$\times$ with no meaningful progress and that is a large warning sign that you need new ideas.}
\end{iband}
\end{frame}

\framekicker{Section 19.3.6 $\cdot$ 2024}
\begin{frame}[shrink=25]{The narrative shift, and what caused it}
{\fontsize{9.4}{12.6}\selectfont\color{fg2}The 2023 bedrock dogma --- \textit{scale is all you need} --- began giving way to \textit{actually, we need on-the-fly recombination}. The ARC Prize results announced in December 2024 were illuminating: \textbf{the leading solutions did not come from scaling existing architectures}.\par}\vskip 4pt
\begin{center}
\includegraphics[width=0.92\linewidth,height=0.60\textheight,keepaspectratio]{figs/ch19-tta.tex.pdf}
\end{center}
\vskip 3pt{\color{ink3}\fontsize{7.4}{9.4}\selectfont Every single top entry in ARC Prize 2024 used test-time adaptation.\par}
{\fontsize{9.4}{12.6}\selectfont\color{fg2}TTA means the system performs \textbf{active reasoning or learning during the test itself}, using the specific problem information provided.\par}\vskip 4pt
\end{frame}

\framekicker{Section 19.3.6 $\cdot$ December 2024}
\begin{frame}[shrink=25]{o3, and the first signs of fluid intelligence}
{\fontsize{9.4}{12.6}\selectfont\color{fg2}Shortly after the competition, OpenAI previewed its \textbf{o3} test-time reasoning model and used ARC-AGI to showcase it.\par}\vskip 4pt
\begin{itable}{>{\raggedright\arraybackslash}p{0.3196\textwidth}>{\raggedright\arraybackslash}p{0.282\textwidth}>{\raggedright\arraybackslash}p{0.3384\textwidth}}
\thead{Compute setting} & \thead{ARC-AGI score} & \thead{Cost per task} \\ \midrule
Moderate & \textbf{76\%} & about \$200 \\
High & \textbf{88\%} --- above the human baseline & over \textbf{\$20,000} \\
\end{itable}
{\fontsize{9.4}{12.6}\selectfont\color{fg2}For the first time, a model showed \textbf{signs of genuine fluid intelligence}. ARC-AGI was one of the only benchmarks at the time giving a clear signal that a paradigm shift was underway.\par}\vskip 4pt
\end{frame}

\framekicker{Section 19.3.7}
\begin{frame}[shrink=25]{Is AGI solved? Not quite --- and here is why}
\begin{minipage}[t]{0.494\textwidth}
\begin{icardwarn}
\cardh{Efficiency is the point}
Tens of thousands of dollars per puzzle. \textbf{Brute-forcing the solution space with enormous compute is a shortcut that makes tasks possible without requiring intelligence.} It felt more like cracking a code with a supercomputer than nimble reasoning.
\end{icardwarn}
\end{minipage}%
\hfill
\begin{minipage}[t]{0.494\textwidth}
\begin{icardwarn}
\cardh{Still stumped by easy tasks}
o3 failed many tasks humans find \textbf{very easy} --- even at the highest compute settings.
\end{icardwarn}
\end{minipage}%
\begin{iband}[signal]
\textbf{The entire point of intelligence is to achieve results with the least resources possible.} In principle you could solve ARC-AGI by enumerating every possible program until one fits the demonstrations. \hilite{That would prove nothing.}
\end{iband}
\end{frame}

\framekicker{Section 19.3.7 $\cdot$ March 2025}
\begin{frame}[shrink=25]{ARC-AGI-2: a benchmark that can tell degrees apart}
{\fontsize{9.4}{12.6}\selectfont\color{fg2}The 2019 version was intended to be \textbf{easy} --- essentially a binary test. Either you have no fluid intelligence and score near zero, or you have some and immediately score very high. There was not much room in between.\par}\vskip 4pt
{\fontsize{9.4}{12.6}\selectfont\color{fg2}ARC-AGI-2 keeps the exact same format with significantly harder content: \textbf{longer reasoning chains, inherently more resistant to exhaustive search}, so that computational efficiency becomes critical to success.\par}\vskip 4pt
\begin{minipage}[t]{0.3253\textwidth}
\istat{instant}{typical human time on an ARC-AGI-1 task}
\end{minipage}%
\hfill
\begin{minipage}[t]{0.3253\textwidth}
\istat{\textasciitilde{}5 min}{average human time on ARC-AGI-2}
\end{minipage}%
\hfill
\begin{minipage}[t]{0.3253\textwidth}
\istat{\textasciitilde{}0\%}{base LLM performance on ARC-AGI-2}
\end{minipage}%
{\fontsize{9.4}{12.6}\selectfont\color{fg2}Even \textbf{o3's scores plummeted into the low double digits} under reasonable compute budgets. The challenge of efficient, human-like fluid intelligence is far from solved.\par}\vskip 4pt
\end{frame}

\coursesection{04}{The missing ingredients: search and symbols}{Two kinds of abstraction, and why any complete mind needs both.}

\framekicker{Section 19.4.1}
\begin{frame}[shrink=25]{Two ways to compare things}
{\fontsize{9.4}{12.6}\selectfont\color{fg2}Abstraction acquisition starts with \textbf{comparing things to one another}. There are two distinct ways to do that, giving rise to two kinds of abstraction and two modes of thinking.\par}\vskip 4pt
\begin{center}
\includegraphics[width=0.92\linewidth,height=0.60\textheight,keepaspectratio]{figs/ch19-two-poles.tex.pdf}
\end{center}
\vskip 3pt{\color{ink3}\fontsize{7.4}{9.4}\selectfont Together, these two poles form the basis for all of our thoughts.\par}
\end{frame}

\framekicker{Section 19.4.1 $\cdot$ figure 19.12}
\begin{frame}[shrink=25]{Value-centric analogy: the beetles in your backyard}
{\fontsize{9.4}{12.6}\selectfont\color{fg2}You notice similarities. Similarity is implicitly a \textbf{smooth, continuous distance function} defining a latent manifold where instances live. Cluster the similar ones and merge them into \textbf{prototypes} capturing shared features.\par}\vskip 4pt
{\fontsize{9.4}{12.6}\selectfont\color{fg2}These prototypes are abstract --- they look like no specific beetle you have seen. Meet a new beetle and you compare it to your handful of prototypes rather than to every beetle in your memory.\par}\vskip 4pt
\begin{iband}[signal]
This is pretty much a description of \textbf{unsupervised machine learning} --- K-means, and in general all of modern machine learning. \hilite{The ConvNet features visualized in chapter 10 were visual prototypes.}
\end{iband}
\end{frame}

\framekicker{Section 19.4.1}
\begin{frame}[shrink=25]{It is also most of what you do without thinking}
\leadpar{\textbf{If you can do a task without thinking about it, you are heavily relying on value-centric analogies.}}
{\fontsize{9.4}{12.6}\selectfont\color{fg2}It underlies pattern recognition, perception, and intuition. Watching a film and subconsciously sorting the characters into \textit{types} is value-centric abstraction.\par}\vskip 4pt
{\fontsize{9.4}{12.6}\selectfont\color{fg2}It is also exactly the kind of analogy-making that enables deep learning models to perform \textbf{local generalization} --- which is the connection this chapter has been building toward.\par}\vskip 4pt
\end{frame}

\framekicker{Section 19.4.1 $\cdot$ figure 19.13}
\begin{frame}[shrink=25]{Program-centric analogy: refactoring}
{\fontsize{9.4}{12.6}\selectfont\color{fg2}In software engineering you write functions and classes that have a lot in common, and start asking: could there be a more abstract function that does the same job twice? An abstract base class both could inherit from?\par}\vskip 4pt
{\fontsize{9.4}{12.6}\selectfont\color{fg2}You are \textbf{not} comparing them by how similar they look, the way you would compare two faces. You are asking whether parts of them have \textbf{exactly the same structure} --- looking for a \textbf{subgraph isomorphism} between programs represented as graphs of operators.\par}\vskip 4pt
\begin{iband}[signal]
This is not exclusive to computer science or mathematics. It underlies reasoning, planning, and the general concept of \textbf{rigour} as opposed to intuition --- any time you think about objects connected by a \hilite{discrete network of relationships}
\end{iband}
\end{frame}

\framekicker{Section 19.4.2 $\cdot$ table 19.1}
\begin{frame}[shrink=25]{The two poles, side by side}
\begin{itable}{>{\raggedright\arraybackslash}p{0.47\textwidth}>{\raggedright\arraybackslash}p{0.47\textwidth}}
\thead{Value-centric abstraction} & \thead{Program-centric abstraction} \\ \midrule
Relates things by \textbf{distance} & Relates things by \textbf{exact structural match} \\
Continuous, grounded in geometry & Discrete, grounded in topology \\
Abstracts by \textit{averaging} instances into prototypes & Abstracts by isolating isomorphic substructures \\
Underlies \textbf{perception and intuition} & Underlies \textbf{reasoning and planning} \\
Immediate, fuzzy, approximate & Slow, exact, rigorous \\
Needs a lot of experience to be reliable & \textbf{Experience efficient} --- can work from two instances \\
\end{itable}
\end{frame}

\framekicker{Section 19.4.2}
\begin{frame}[shrink=25]{Nothing you do uses only one of them}
{\fontsize{9.4}{12.6}\selectfont\color{fg2}You would be hard pressed to find a task involving only one pole. Even \textit{pure perception} --- recognizing objects in a scene --- involves implicit reasoning about the relationships between them.\par}\vskip 4pt
{\fontsize{9.4}{12.6}\selectfont\color{fg2}And even \textit{pure reasoning} --- finding the proof of a theorem --- involves a good deal of intuition. \textbf{When a mathematician puts pen to paper they already have a fuzzy vision of the direction they are going.} The discrete steps are guided by high-level intuition.\par}\vskip 4pt
\begin{iband}[signal]
These two poles are complementary, and \textbf{it is their interleaving that enables extreme generalization}. No mind could be complete without both.
\end{iband}
\end{frame}

\framekicker{Section 19.4.3}
\begin{frame}[shrink=25]{Sorting five numbers, and classifying MNIST}
\begin{minipage}[t]{0.494\textwidth}
\begin{icardbad}
\cardh{Sorting, with deep learning}
Possible with the right architecture, and \textbf{an exercise in frustration}. Massive data, occasional mistakes on new numbers, and a complete retrain to go from lists of 5 to lists of 10. A Python sort is a few lines and works on \textbf{any} list.
\end{icardbad}
\end{minipage}%
\hfill
\begin{minipage}[t]{0.494\textwidth}
\begin{icardbad}
\cardh{MNIST, with discrete reasoning}
Hand-code closed-loop counts, centres of mass. After \textbf{thousands of lines} you might reach 90\% test accuracy. Fitting a parametric model is simpler, uses the data better, and is far more robust.
\end{icardbad}
\end{minipage}%
\begin{iband}[signal]
So it is unlikely anyone will reduce reasoning to manifold interpolation, or perception to discrete reasoning. \textbf{The way forward is a unified framework incorporating both.}
\end{iband}
\end{frame}

\framekicker{Section 19.4.4 $\cdot$ figure 19.14}
\begin{frame}[shrink=25]{Program synthesis, the missing piece}
\begin{center}
\includegraphics[width=0.92\linewidth,height=0.60\textheight,keepaspectratio]{figs/ch19-program-synth.tex.pdf}
\end{center}
\vskip 3pt{\color{ink3}\fontsize{7.4}{9.4}\selectfont Highly reminiscent of machine learning --- given input-output pairs, find a program that generalizes.\par}
{\fontsize{9.4}{12.6}\selectfont\color{fg2}The difference: instead of learning \textbf{parameter values} inside a hardcoded program (a neural network), we \textbf{generate source code} through a discrete search process.\par}\vskip 4pt
\end{frame}

\framekicker{Section 19.4.4 $\cdot$ table 19.2}
\begin{frame}[shrink=25]{Machine learning against program synthesis}
\begin{itable}{>{\raggedright\arraybackslash}p{0.188\textwidth}>{\raggedright\arraybackslash}p{0.376\textwidth}>{\raggedright\arraybackslash}p{0.376\textwidth}}
\thead{} & \thead{Machine learning} & \thead{Program synthesis} \\ \midrule
Model & A differentiable parametric function & A graph of operators from a programming language \\
Engine & \textbf{Gradient descent} & \textbf{Discrete search} (e.g. genetic search) \\
Data & Requires a lot to be reliable & \textbf{Data efficient} --- works from a couple of examples \\
\end{itable}
{\fontsize{9.4}{12.6}\selectfont\color{fg2}Until 2024, AI systems capable of genuine discrete reasoning were \textbf{all hardcoded by human programmers}. In the test-time adaptation era that is finally changing --- and program synthesis is the branch to watch.\par}\vskip 4pt
\end{frame}

\framekicker{Section 19.4.5}
\begin{frame}[shrink=25]{Hybrids already exist, and they are hand-built}
\begin{minipage}[t]{0.494\textwidth}
\begin{icardaccent}
\cardh{AlphaGo}
\textbf{Most of the intelligence on display is hardcoded by human programmers} --- Monte Carlo Tree Search. Learning from data happens only in specialized submodules: value and policy networks.
\end{icardaccent}
\end{minipage}%
\hfill
\begin{minipage}[t]{0.494\textwidth}
\begin{icardaccent}
\cardh{A Waymo self-driving car}
Maintains a literal \textbf{3D model of the world}, full of assumptions hardcoded by engineers, constantly updated by deep learning perception modules --- powered by Keras.
\end{icardaccent}
\end{minipage}%
{\fontsize{9.4}{12.6}\selectfont\color{fg2}In both, the combination of human-created discrete programs and learned continuous models unlocks performance \textbf{impossible with either alone}. Today the discrete parts are painstakingly hand-coded. In the future such systems may be fully learned.\par}\vskip 4pt
\end{frame}

\framekicker{Section 19.4.5 $\cdot$ a thought experiment}
\begin{frame}[shrink=25]{An RNN has one hardcoded for loop. What if it had more?}
{\fontsize{9.4}{12.6}\selectfont\color{fg2}RNNs have slightly fewer limitations than feedforward networks because they are more than a geometric transformation: they are \textbf{a geometric transformation applied repeatedly inside a for loop}. That loop is hardcoded by human developers --- a built-in assumption.\par}\vskip 4pt
{\fontsize{9.4}{12.6}\selectfont\color{fg2}Now imagine a network augmented not with a single hardcoded loop and continuous-space memory, but with a large set of \textbf{programming primitives} it is free to manipulate: \inlinecode{if} branches, \inlinecode{while} statements, variable creation, disk storage, sorting operators, lists, graphs, hash tables.\par}\vskip 4pt
\begin{iband}[signal]
The space of programs such a network could represent would be far broader --- and \textbf{some would generalize far better}. Such programs will not be differentiable end to end, so they must be generated by \hilite{a combination of discrete program search and gradient descent}
\end{iband}
\end{frame}

\framekicker{Section 19.4.5 $\cdot$ the obstacle}
\begin{frame}[shrink=25]{Combinatorial explosion}
{\fontsize{9.4}{12.6}\selectfont\color{fg2}Program synthesis today is \textbf{tremendously inefficient}. To caricature: it tries every possible program in a search space until one matches the specification.\par}\vskip 4pt
{\fontsize{9.4}{12.6}\selectfont\color{fg2}As specification complexity rises, or the vocabulary of primitives expands, the set of possible programs grows \textbf{much faster than merely exponentially}. So today, program synthesis generates only very short programs. \textit{You are not going to be generating a new OS for your computer anytime soon.}\par}\vskip 4pt
\begin{iband}[signal]
The fix is to bring it closer to how humans write software. When you open your editor you are \textbf{not thinking about every possible program} --- you have a handful of approaches in mind, cut down by understanding and past experience.
\end{iband}
\end{frame}

\framekicker{Section 19.4.5 $\cdot$ the resolution}
\begin{frame}[shrink=25]{Deep learning can guide the search}
{\fontsize{9.4}{12.6}\selectfont\color{fg2}Each specific program you want to generate is a fundamentally discrete object doing non-interpolative manipulation. But evidence so far indicates that \textbf{the space of all useful programs may look a lot like a continuous manifold}.\par}\vskip 4pt
{\fontsize{9.4}{12.6}\selectfont\color{fg2}A model trained on millions of successful program-generation episodes might develop \textbf{solid intuition about the path through program space} --- just as an engineer has immediate intuition about the architecture of the script they are about to write.\par}\vskip 4pt
\begin{iband}[signal]
Human reasoning is itself heavily guided by value-centric abstraction --- by pattern recognition and intuition. \textbf{The same should be true of program synthesis.} Expect increasing research interest over the next 10 to 20 years.
\end{iband}
\end{frame}

\framekicker{Section 19.4.6}
\begin{frame}[shrink=25]{What LLMs are missing, stated precisely}
{\fontsize{9.4}{12.6}\selectfont\color{fg2}Foundation models bring us closer to systems with enormous acquired knowledge and skills. But LLMs are missing \textbf{recombination}: they fetch and reapply memorized functions well, and cannot recombine them on the fly into new programs.\par}\vskip 4pt
{\fontsize{9.4}{12.6}\selectfont\color{fg2}They are in fact \textbf{entirely incapable of function composition}, as investigated by Dziri et al. And the functions they learn are not sufficiently abstract or modular to be recombined in the first place.\par}\vskip 4pt
\begin{iband}[amber]
Remember the low accuracy on adding large integers? \hilite{You would not want to build your next codebase on top of such brittle functions.}
\end{iband}
\end{frame}

\framekicker{Section 19.4.6 $\cdot$ figure 19.16}
\begin{frame}[shrink=25]{A global library, and a perpetual meta-learner}
\begin{center}
\includegraphics[width=0.92\linewidth,height=0.60\textheight,keepaspectratio]{figs/ch19-hybrid.tex.pdf}
\end{center}
\vskip 3pt{\color{ink3}\fontsize{7.4}{9.4}\selectfont Subroutines may be geometric (pretrained deep learning modules) or algorithmic (closer to the libraries engineers use today).\par}
{\fontsize{9.4}{12.6}\selectfont\color{fg2}Think of software development: once an engineer solves HTTP queries in Python, they package it as a reusable library accessible to anyone on the planet. \textbf{Any single problem encountered would only need to be solved once} --- making such systems constantly self-improving.\par}\vskip 4pt
\end{frame}

\framekicker{Section 19.4.7}
\begin{frame}[shrink=25]{The long-term vision, in four points}
\begin{isteps}
\item \textbf{Models will be more like programs}, going far beyond continuous geometric transformations --- closer to the abstract mental models humans maintain, and capable of stronger generalization.
\item \textbf{They will blend algorithmic modules} providing formal reasoning, search, and abstraction \textbf{with geometric modules} providing intuition and pattern recognition. AlphaGo and self-driving cars are an early, hand-built example.
\item \textbf{They will be grown automatically} rather than hardcoded, from modular parts in a global library evolved across thousands of previous tasks --- frequent problem-solving patterns turned into reusable subroutines.
\item \textbf{The search over combinations will be discrete} --- program synthesis --- but heavily guided by a form of program-space intuition provided by deep learning.
\end{isteps}
\end{frame}

\framekicker{What to carry into practice}
\begin{frame}[shrink=25]{Four things this chapter should change about your work}
\begin{minipage}[t]{0.494\textwidth}
\begin{icardaccent}
\cardh{Ask about familiarity, not difficulty}
When scoping an LLM deployment, the question is whether the task resembles something in web text --- \textbf{not whether it is hard}.
\end{icardaccent}
\end{minipage}%
\hfill
\begin{minipage}[t]{0.494\textwidth}
\begin{icardaccent}
\cardh{Test with perturbations}
Change names, numbers, and phrasing that do not change the problem. \textbf{If the answer changes, you have measured memorization.}
\end{icardaccent}
\end{minipage}%


\vskip 4pt

\begin{minipage}[t]{0.494\textwidth}
\begin{icardaccent}
\cardh{Choose the metric carefully}
The shortcut rule is not a research curiosity. Optimize one number and you will get it, \textbf{at the expense of everything else} --- as Netflix did.
\end{icardaccent}
\end{minipage}%
\hfill
\begin{minipage}[t]{0.494\textwidth}
\begin{icardaccent}
\cardh{Expect hybrids}
The systems that work in the real world today combine learned perception with hand-built discrete logic. \textbf{Design for that combination}, not for an end-to-end network.
\end{icardaccent}
\end{minipage}%
\end{frame}

\framekicker{Summary}
\begin{frame}[shrink=25]{What has to survive this chapter}
\begin{isteps}
\item \textbf{A deep learning model is a static interpolative database.} The only operation it supports is retrieval.
\item \textbf{Success tracks familiarity, not complexity.} A trivially easy but novel problem defeats every current model.
\item \textbf{Prompt engineering is search}, not communication. Its necessity is evidence against understanding.
\item \textbf{Scale has not moved any of these limits} in seven years and four orders of magnitude, because the model is unchanged.
\end{isteps}
\end{frame}

\framekicker{Summary}
\begin{frame}[shrink=25]{What has to survive this chapter (2 of 2)}
\begin{isteps}
\item \textbf{Local, broad, and extreme generalization} are different things. Deep learning does the first, and only for anticipated factors of variation.
\item \textbf{Fixing the task removes the need for intelligence} --- and an unintelligent solution is always the easier one, so it is the one you get.
\item \textbf{Intelligence is an efficiency ratio} between information available and future operating area --- which is what ARC-AGI measures.
\item \textbf{Test-time adaptation was the 2024 shift}: every top ARC Prize entry used it, and o3 crossed the human baseline --- at \$20,000 a task.
\item \textbf{Two poles of abstraction}: value-centric for perception, program-centric for reasoning. Deep learning has one of them.
\item \textbf{The way forward is a blend} --- program synthesis for discrete structure, deep learning to guide the search.
\end{isteps}
\vskip 4pt
\reslink{PAPER}{Chollet, On the Measure of Intelligence (2019)}{https://arxiv.org/abs/1911.01547}
\reslink{SITE}{ARC Prize}{https://arcprize.org/}
\reslink{NEXT}{Chapter 20 --- Conclusions}{../ch20/index.html}
\vskip 2pt
\end{frame}

\end{document}
